% Radial-identity repair / execution appendix for DASHIAtomicPeriodicTable369Formalism.tex
% Repo-native owners:
%   DASHI/Physics/Chemistry/AtomicPeriodicTable369ArchiveHartreeRadialIdentityDefectExact.agda
%   DASHI/Physics/Chemistry/AtomicPeriodicTable369RadialIdentityRepairExecutionExact.agda
%   DASHI/Physics/Chemistry/AtomicPeriodicTable369PhysicalValidation.agda

\section{Radial-state identity repair and execution receipt}

The historical \emph{DASHI Atom} Hartree code-shaped object accepted subshell
triples $(n,\ell,\mathrm{occ})$, but passed only $\ell$ into its radial
eigenproblem and always requested the lowest eigenpair in that $\ell$ sector.
Consequently $1s$, $2s$, $3s$, \ldots{} collapsed to the same radial target.
This defect is preserved append-only in
\texttt{AtomicPeriodicTable369ArchiveHartreeRadialIdentityDefectExact.agda}; the
later repair does not rewrite the historical source state.

\subsection{Corrected radial selector}

For an analytically controlled hydrogenic central potential, the radial state
index associated with a bound label $(n,\ell)$ is
\[
  r(n,\ell)=n-\ell-1.
\]
The repaired producer
\[
  \texttt{scripts/atomic\_radial\_identity\_repair.py}
\]
therefore diagonalizes each complete $\ell$ sector far enough to contain every
required radial index and selects eigenpair $r(n,\ell)$.  Solving the sector as
one eigenproblem also gives a direct within-$\ell$ orthogonality check rather
than solving the same lowest state repeatedly.

The committed execution receipt is
\[
\texttt{Artifacts/atomic-periodic-table-369/radial-identity-repair-receipt.json}
\]
with producer SHA--256
\[
\texttt{d4df9c2ef1e2292cdf216d7bf5da4aab78189b47cff02dc83bb28d308d240b9e}.
\]
It records Python 3.13.5, NumPy 2.3.5 and SciPy 1.17.0, $4000$ radial grid
points on $r\in(0,80]$ Bohr, and an energy tolerance of $5\times 10^{-4}$
Hartree.

\subsection{Executed hydrogenic regression}

For $Z=1$ and $V_H=0$, the numerical producer checks the analytic energy law
\[
  E_n=-\frac{1}{2n^2}\ \text{Hartree}
\]
and radial node count $n-\ell-1$.  The retained values are:

\begin{center}
\begin{tabular}{lrrrr}
\toprule
state & radial index & nodes & computed $E$ & analytic $E$\\
\midrule
$1s$ & 0 & 0 & $-0.4999500349782737$ & $-0.5$\\
$2s$ & 1 & 1 & $-0.12499687671803333$ & $-0.125$\\
$3s$ & 2 & 2 & $-0.05555493859367163$ & $-1/18$\\
$2p$ & 0 & 0 & $-0.12500104120998212$ & $-0.125$\\
$3p$ & 1 & 1 & $-0.055556035451670566$ & $-1/18$\\
$3d$ & 0 & 0 & $-0.055555596687603834$ & $-1/18$\\
\bottomrule
\end{tabular}
\end{center}

All energy, node-count and within-$\ell$ orthogonality checks passed.  The
largest retained Gram-matrix deviation was below $7\times10^{-16}$.

\subsection{Attribution snowball}

The identifier fields remain source-role separated:

\begin{itemize}
\item Schr\"odinger, \emph{Quantisierung als Eigenwertproblem} (1926), DOI
  \texttt{10.1002/andp.19263840404}, is retained as a primary historical
  eigenvalue-problem source.  Semantic coordinates remain atom Q9121 and
  quantum mechanics Q944, with Dewey topic coordinates 539.7 and 530.12.
\item NIST Atomic Spectra Database, Standard Reference Database 78, DOI
  \texttt{10.18434/T4W30F}, is the authoritative downstream comparison source
  for atomic energy levels, ground states and ionization energies.  Ionization
  energy remains Q483769.  Database presence is not agreement.
\item OEIS A167268, A093907 and A018227 remain upstream structural-selector,
  period-length and closure coordinates.  They do not validate the radial
  eigenproblem or quantitative atomic energies.
\end{itemize}

\subsection{Promotion boundary}

The newly paid statement is therefore narrow but executable:
\[
\boxed{
(n,\ell)\mapsto\text{distinct radial eigenstate}
\quad\text{is implemented and numerically regressed in the hydrogenic limit.}
}
\]
It does \emph{not} imply
\[
\text{SCF convergence},\quad
\text{correct many-electron ground states},\quad
\text{calibrated ionization energies},\quad
\text{nuclear stability},\quad
\text{empirical periodic-table recovery}.
\]
The next same-object producer must carry the repaired state identities through
self-consistent iterations, specify self-interaction/exchange treatment, emit
convergence diagnostics and total energies for matched neutral/cation objects,
and only then compare
\[
  IE_1(Z)=E(Z-1)-E(Z)
\]
against NIST ASD with uncertainty-aware residuals.
